\chapter{Deliverable}
\label{chap:deliverable}

\noindent
This chapter summarizes the tangible outputs produced by the project and the test/evaluation protocol used to assess them. It focuses on what is delivered (software and artifacts) and how performance is measured for RT regression and correctness classification.
\vspace{1em}

\section{Software Solution}
\label{sec:software-solution}
\begin{itemize}
	\item \textbf{Project repository:} Contains preprocessing scripts and modeling code. [Link: to be added]
	\item \textbf{Reproducible pipeline artifacts:} dataset construction routines for CCD-ITI windows, preprocessing configurations, feature extraction scripts, and run logging.
	\item \textbf{Figures:} verification plots (e.g., SNR spectra) and evaluation plots used in the report; architecture diagram(s) documenting the workflow.
	\item \textbf{Interface documentation/prototype:} UI screenshots and descriptions supporting the experiment workflow.
\end{itemize}

\section{Test Report}
\label{sec:test-report}
\begin{itemize}
	\item \textbf{CCD-ITI approach:} Train on CCD-ITI segments and evaluate under a defined split strategy (preferably subject-independent).
	\item \textbf{Regression metrics (RT):} MAE, RMSE, R², and Pearson correlation.
	\item \textbf{Classification metrics (correctness):} accuracy, balanced accuracy, sensitivity, and specificity.
	\item \textbf{Verification:} Use SNR spectra (and related plots) to confirm steady-state characteristics before interpreting model outcomes.
\end{itemize}

\subsection{Classification Runs}
\label{subsec:classification-failure-summary}

\begin{table}[H]
\centering
\begin{tabular}{|l|c|c|c|c|}
\hline
\textbf{Balance Strategy} & \textbf{Acc} & \textbf{Bal Acc} & \textbf{Sens} & \textbf{Spec} \\
\hline
None & 0.729 & 0.501 & 0.993 & 0.009 \\
\hline
Class Weight & 0.604 & 0.544 & 0.674 & 0.414 \\
\hline
Weighted Sampler & 0.629 & 0.543 & 0.729 & 0.358 \\
\hline
Undersample & 0.652 & 0.537 & 0.787 & 0.287 \\
\hline
\end{tabular}
\caption{Selected classification runs from W\&B export.}
\label{tab:classification-failure-summary}
\end{table}

\noindent
\textbf{Interpretation:} None shows majority-class collapse (almost all predictions biased to the positive class). Class Weight achieves better balance but still exhibits class-bias and limited generalization. Weighted Sampler helps recall but sacrifices specificity. Undersample improves minority exposure, though instability and false-positive pressure remain.

\subsection{Regression Runs}
\label{subsec:regression-failure-summary}

\begin{table}[H]
\centering
\begin{tabular}{|l|c|c|c|c|}
\hline
\textbf{Loss} & \textbf{MAE} & \textbf{RMSE} & \textbf{R²} & \textbf{Pearson} \\
\hline
MAE & 0.2876 & 0.3743 & 0.0063 & 0.0840 \\
\hline
MSE & 0.2883 & 0.3749 & 0.0030 & 0.0562 \\
\hline
Huber & 0.2874 & 0.3740 & 0.0079 & 0.0903 \\
\hline
\end{tabular}
\caption{Regression runs from W\&B export.}
\label{tab:regression-failure-summary}
\end{table}

\noindent
\textbf{Interpretation:} MAE shows low explanatory power with predictions remaining close to average RT. MSE exhibits similar stagnation, suggesting that loss choice alone does not improve structure learning. Huber performs slightly better in correlation but still shows near-zero R² and weak trial-level fit.

\noindent
Overall, no single run is presented as ``best.'' The current test report is used to identify concrete failure modes that guide the next iteration.
